% This file was converted to LaTeX by Writer2LaTeX ver. 1.9.9
% see http://writer2latex.sourceforge.net for more info
\documentclass{article}
\usepackage{calc,amsmath,amssymb,amsfonts}
\usepackage[LGR,T1]{fontenc}
\usepackage[greek,italian]{babel}
\usepackage[style=numeric,backend=biber]{biblatex}
\author{HP}
\date{2026-07-31}
\begin{document}
\section{Introduzione}
Ci sono ricordi che il tempo sfuma lentamente, fino a renderli quasi irriconoscibili.

Altri, invece, rimangono sorprendentemente nitidi. Non sapresti dire quale giorno fosse, né ricordare con precisione il
mese o l'ora. Eppure, se chiudi gli occhi, riesci ancora a rivedere la stanza, la luce che entrava dalla finestra e
perfino la sensazione provata in quell'istante.

Il mio primo incontro con un videogioco appartiene a questa seconda categoria.

Era il 1982.

Frequentavo la scuola media quando un pomeriggio andai a trovare un mio compagno di classe, Emanuele. Era una visita
come tante altre. Avremmo trascorso qualche ora insieme, parlando di scuola, giocando o semplicemente facendo quello
che fanno due ragazzi della stessa età. Non avevo la minima idea che quel pomeriggio avrebbe lasciato un segno
destinato a rimanere con me per tutta la vita.

Entrai nella sua cameretta. Ricordo una scrivania appoggiata contro la parete, un televisore acceso e lui che, con
naturalezza, prese in mano un joystick. Sullo schermo comparvero due barre bianche e un piccolo quadrato luminoso che
continuava a rimbalzare da una parte all'altra.

Era Pong.

Oggi potrebbe sembrare poco più di una curiosità storica. Ma allora era qualcosa di completamente diverso. Per me non
era soltanto un videogioco.

Era una domanda.

Guardavo quella piccola pallina cambiare direzione, rimbalzare, assegnare un punto, ricominciare la partita. Più
osservavo lo schermo, meno mi interessava giocare. Continuavo a chiedermi come fosse possibile. Come faceva quel
televisore a sapere dove si trovasse la pallina? Come decideva quando invertire la direzione? Come riconosceva una
collisione? Come stabiliva il punteggio?

Quando tornai a casa non avevo ancora trovato una risposta.

Ma avevo trovato la domanda che avrebbe accompagnato gran parte della mia vita.

Come fa?

Negli anni successivi quella domanda non smise più di presentarsi. Ogni volta che vedevo un computer, ogni volta che
osservavo un videogioco, ogni volta che sentivo parlare di programmazione, tornava puntualmente a bussare nella mia
mente.

Non sapevo ancora che la risposta sarebbe arrivata l'anno successivo, grazie a un computer che portava un numero
apparentemente insignificante.

Sessantaquattro.


\bigskip

Quel computer entrò nella mia vita quando ero ancora un ragazzo e, quasi senza accorgermene, cambiò il modo in cui
osservavo il mondo. Non mi limitavo più a utilizzare ciò che avevo davanti. Cominciai a chiedermi come fosse stato
costruito. Ogni comando imparato, ogni programma scritto, ogni errore corretto sembravano aggiungere un piccolo
tassello a una risposta che appariva ancora lontana.

Come molti ragazzi della mia generazione, anch'io iniziai presto a coltivare un sogno.

Non desideravo semplicemente imparare il BASIC.

Non desideravo soltanto giocare.

Volevo creare un videogioco vero, uno di quelli che riuscivano a raccontare una storia, a dare vita a un personaggio, a
trasformare poche decine di kilobyte in un mondo capace di far immaginare molto più di quanto mostrasse sullo schermo.

Quel videogioco, però, non arrivò mai.

La vita prese altre strade. Gli studi, il lavoro, le responsabilità e i progetti professionali finirono, poco alla
volta, per occupare tutto lo spazio. Il Commodore 64 venne riposto in una scatola insieme ai suoi cavi, al Datasette,
ai manuali e ai joystick consumati dall'uso.

Pensavo che quella storia fosse finita. Mi sbagliavo.

Molti anni dopo mi capitò di ritrovare quella scatola. Non ricordo nemmeno cosa stessi cercando. Forse una vecchia
valigia. Forse alcuni album fotografici. Ricordo soltanto la polvere, il cartone ormai ingiallito e la sensazione di
avere improvvisamente aperto una finestra sul passato.

Dentro non c'era soltanto un vecchio computer.

C'era una parte della mia storia.

Presi il Commodore 64 tra le mani e rimasi per qualche minuto in silenzio.

Poi sorrisi.

$\text{\textgreek{«}}$Ci ho messo quarant'anni.$\text{\textgreek{»}}$

Non quarant'anni per ritrovare quel computer.

Quarant'anni per capire che esisteva ancora una promessa rimasta incompiuta.

Il bambino che aveva desiderato con tutte le sue forze quella macchina aveva mantenuto quasi tutte le promesse fatte a
sé stesso. Grazie a quel computer aveva scoperto la programmazione, scelto il proprio percorso di studi e trasformato
una curiosità in una professione.

Ma una promessa era rimasta indietro.

Realizzare quel videogioco che aveva immaginato tante volte e che non aveva mai avuto il coraggio, o forse semplicemente
l'occasione, di cominciare.

È da quel momento che nasce questo libro.

Non dall'idea di scrivere un manuale sul BASIC del Commodore 64.

Nemmeno dal desiderio di raccogliere alcuni ricordi personali.

Entrambe queste componenti sono presenti, ma nessuna rappresenta il vero motivo per cui queste pagine esistono.

Questo libro nasce dal desiderio di chiudere un cerchio rimasto aperto per oltre quarant'anni.

Ma non voglio farlo da solo.

Vorrei che il lettore percorresse questo cammino insieme a me.

Non importa se conserva ancora il proprio Commodore 64 perfettamente funzionante, se lo ha ritrovato dopo decenni in una
cantina o se lo conoscerà soltanto attraverso un emulatore. Non importa nemmeno se ha vissuto gli anni Ottanta o se
appartiene a una generazione che li ha conosciuti soltanto attraverso i racconti.

Quello che conta è la curiosità.

La stessa curiosità che nasce quando ci si trova davanti a qualcosa che sembra quasi magico e si decide di non
accontentarsi di usarlo, ma di capire come sia stato costruito.

Per questo motivo non affronteremo questo percorso come un manuale tradizionale.

Non sarà la teoria a guidare il progetto. Sarà il progetto stesso a indicarci, passo dopo passo, quali domande porci,
quali strumenti imparare e perché ci servono. Costruiremo un videogioco così come avrebbe potuto nascere negli anni
Ottanta: un problema alla volta, una scelta alla volta, un errore alla volta. Ogni nuova conoscenza nascerà da una
necessità concreta. Ogni difficoltà diventerà l'occasione per imparare qualcosa di nuovo. Ogni soluzione farà crescere
il progetto.

Il lettore non assisterà semplicemente allo sviluppo del programma.

Ne farà parte.

Vedrà il videogioco nascere, cambiare, sbagliare e migliorare. Scoprirà che i limiti del Commodore 64 non erano soltanto
ostacoli, ma strumenti che insegnavano a ragionare. Capirà che programmare significa molto più che impartire istruzioni
a una macchina: significa imparare a osservare un problema, formulare un'ipotesi e verificare se quell'idea sia davvero
in grado di funzionare.

Accanto al percorso tecnico ne scorrerà un altro, più personale.

È la storia delle persone che hanno alimentato quella curiosità, degli incontri che hanno lasciato un segno, delle
occasioni che hanno cambiato una direzione e delle esperienze che hanno trasformato un bambino incuriosito da una
pallina che rimbalzava sullo schermo in un programmatore.

Non so quale sarà il risultato di questo viaggio per ciascun lettore.

Forse qualcuno ritroverà emozioni che credeva dimenticate. Qualcun altro incontrerà per la prima volta una macchina
appartenente a un'altra epoca e scoprirà quanto possa ancora insegnare.

Mi auguro soltanto che, arrivato all'ultima pagina, il lettore non abbia semplicemente imparato qualcosa in più sul
Commodore 64.

Spero che abbia avuto la sensazione di aver partecipato davvero alla nascita di un videogioco.

E, soprattutto, che abbia scoperto il piacere di guardare un computer in modo diverso.

Perché programmare, in fondo, significa proprio questo: smettere di domandarsi soltanto che cosa una macchina sappia
fare e cominciare a chiedersi come sia possibile insegnarglielo.

È da domande come questa che nascono i programmatori.

Ed è da una di quelle domande che, molti anni fa, è cominciato anche il mio viaggio.


\bigskip
\end{document}
